\section{Introduction}

Throughout the paper, all graphs are finite, simple, and undirected,
and we use standard terminology and notation
\cite{bondy_graph_2008,diestel_graph_2017}.
For a vertex $v$ of a graph $G$, let $N_G(v)$ and $d_G(v)$ denote its
neighborhood and degree, respectively.
For a positive integer $r$, an \emph{$r$-dynamic coloring} of a graph
$G$ is a proper vertex coloring $c$ such that
\[
|c(N_G(v))|\ge \min\{r,d_G(v)\}
\]
for every $v\in V(G)$. The minimum number of colors in such a coloring
is the \emph{$r$-dynamic chromatic number} of $G$, denoted by
$\chi_r(G)$. Dynamic coloring, corresponding to $r=2$, was introduced
by Montgomery~\cite{montgomery2001dynamic}. Higher-order versions have
subsequently been studied under the names conditional coloring,
$r$-hued coloring, and $r$-dynamic coloring; see
\cite{CHEN202224,LAI20061997}.

Dynamic coloring provides a natural interpolation between ordinary
coloring and coloring the square of a graph. Recall that $G^2$ has
vertex set $V(G)$, with two distinct vertices adjacent whenever their
distance in $G$ is at most two. For every graph $G$,
\begin{equation}\label{eq:absolute-bounds}
\chi(G)=\chi_1(G)
\le \chi_2(G)
\le\cdots\le
\chi_{\Delta(G)}(G)
=\chi(G^2).
\end{equation}
In particular, if $G$ is $d$-regular, then
\begin{equation}\label{eq:regular-square}
\chi_d(G)=\chi(G^2).
\end{equation}
Thus, on regular graphs, the last member of the dynamic-coloring
hierarchy is precisely the classical square-coloring problem.

Both dynamic coloring and coloring graph squares have received
considerable attention for graphs of bounded degree. We write
$\text{ch}(G)$ for the list chromatic number of $G$.
Cranston and Kim~\cite{CranstonKim2008} proved that every connected
subcubic graph other than the Petersen graph satisfies
$
\text{ch}(G^2)\le8,
$
and that the bound is best possible. For planar cubic graphs,
Thomassen~\cite{Thomassen2018} proved the sharp bound
$
\chi(G^2)\le7.
$
Square coloring has also been studied under structural restrictions.
In particular, Cames van Batenburg and Kang
\cite{CamesVanBatenburgKang2019} proved that every claw-free graph
with clique number three satisfies
$
\chi(G^2)\le10,
$
and further bounds for squares of claw-free graphs were obtained in
\cite{DeJoannisKangPastor2019}.

The regular setting is also natural from the dynamic-coloring
perspective. Montgomery conjectured that every regular graph satisfies
$
\chi_2(G)\le\chi(G)+2.
$
Bowler et al.~\cite{BowlerEtAl2017} disproved this conjecture by showing
that for every integer $k\ge2$ there exists a regular graph $G$ with
$
\chi(G)=k
\qquad\text{and}\qquad
\chi_2(G)=2k.
$
Thus, regularity alone does not force the dynamic chromatic number to
remain close to the ordinary chromatic number, motivating the study of
structured classes of regular graphs.

Our first main result concerns claw-free cubic graphs. For such graphs,
the $3$-dynamic coloring problem is exactly the square-coloring problem
by~\eqref{eq:regular-square}. Li and Lai~\cite{LiLai2017} proved that
every claw-free graph satisfies
$
\chi_3(G)\le \max\{\chi(G)+3,7\},
$
which yields the upper bound $7$ for claw-free cubic graphs. We improve
this bound to $6$, and the improvement is best possible.

\begin{itemize}
    \item 
If $G$ is a claw-free cubic graph, then
$
\chi(G^2)=\chi_3(G)\le6.
$
Moreover, the bound is sharp. (Theorem \ref{thm:clawfreeleq6}).
\end{itemize}
%\begin{theorem}\label{thm:clawfree-cubic}
%If $G$ is a claw-free cubic graph, then
%$
%\chi(G^2)=\chi_3(G)\le6.
%$
%Moreover, the bound is sharp.
%\end{theorem}

Indeed, the triangular prism $C_3\square K_2$ has diameter two, and
hence its square is $K_6$.
The proof of Theorem \ref{thm:clawfreeleq6} combines the structure of claw-free cubic
graphs with incidence colorings of subcubic multigraphs. We then exploit
the structure more strongly in the Hamiltonian case and reduce the
bound by one, apart from four explicit graphs (to be defined later).

%\begin{theorem}\label{thm:clawfree-hamiltonian}
%Let $G$ be a Hamiltonian claw-free cubic graph. If $ G\notin \{C_3\square K_2,X_{12},X_{16}^{1},X_{16}^{2}\}, $ then
%$
%\chi(G^2)=\chi_3(G)\le5.
%$
%\end{theorem}

\begin{itemize}
    \item If $G$ is a Hamiltonian claw-free cubic graph then
$
\chi(G^2)=\chi_3(G)\le5$, except
for four particular graphs. (Theorem \ref{thm:ham-clawfree}).
\end{itemize}

For completeness, observe the case $r=2$ is immediate for claw-free cubic graphs. The neighborhood of every vertex contains an edge, and hence
every proper coloring is automatically $2$-dynamic. Consequently,
$
\chi_2(G)=\chi(G).
$

Our second direction is to determine the complete dynamic-coloring
hierarchy for a natural family of $4$-regular graphs. For positive
integers $p,n_1,\ldots,n_k$, the circulant graph
$C_p(n_1,\ldots,n_k)$ has vertex set $\mathbb Z_p$, with $i$ adjacent
to $i\pm n_j$ for every $j$. We focus on $C_p(1,3)$ and determine
$\chi_r(C_p(1,3))$ for every $r\in\{2,3,4\}$.

The following hold (Theorems
\ref{thm:circulant-r2},     
\ref{thm:circulant-r3}, and \ref{thm:circulant-r4}):
\begin{itemize}
 \item If $p\ge6$, then
\[
\chi_2(C_p(1,3))=
\begin{cases}
3,&9\mid p,\\
4,&\text{otherwise}.
\end{cases}
\]

\item If $p\ge12$, then
\[
\chi_3(C_p(1,3))=
\begin{cases}
4,&12\mid p,\\
5,&\text{otherwise}.
\end{cases}
\]

\item If $p\ge12$, then
\[
\chi_4(C_p(1,3))=
\begin{cases}
5,&5\mid p,\\
8,&p=16,\\
7,&p\in\{13,14,19,21,26,31\},\\
6,&\text{otherwise}.
\end{cases}
\]
\end{itemize}

Since $C_p(1,3)$ is $4$-regular for $p\ge12$, the last part also
determines the chromatic number of its square. In fact,
$
C^2_p(1,3)=C_p(1,2,3,4,6).
$

%Thus, the two parts of the paper address complementary aspects of dynamic coloring in low-degree regular graphs: sharp square-coloring bounds for a structural class of cubic graphs, and an exact determination of the entire dynamic-coloring hierarchy for the circulant family $C_p(1,3)$.